% !TEX root = ../1Main.tex
% Einleitung - Template/Beispielkapitel

\chapter{Einleitung}
\label{chap:einleitung}

% ===================================================================
% ANLEITUNG: Dieses Kapitel dient als Vorlage für Ihre Einleitung.
% Passen Sie die Struktur und den Inhalt an Ihre Arbeit an.
% ===================================================================

Dieses Kapitel zeigt die grundlegende Struktur einer Einleitung und enthält Beispiele für wichtige Formatierungselemente.

\section{Motivation und Problemstellung}
\label{sec:motivation}

In diesem Abschnitt beschreiben Sie die \textbf{Motivation} für Ihre Arbeit und die zu lösende \textit{Problemstellung}. Nutzen Sie \textbf{Fettdruck} für Hervorhebungen und \textit{Kursivschrift} für Begriffe.

Beispiel für eine Aufzählung:
\begin{itemize}
    \item Erster wichtiger Punkt
    \item Zweiter wichtiger Punkt
    \item Dritter wichtiger Punkt
\end{itemize}

Nummerierte Listen sind ebenfalls möglich:
\begin{enumerate}
    \item Erste Anforderung
    \item Zweite Anforderung
    \item Dritte Anforderung
\end{enumerate}

\section{Zielsetzung}
\label{sec:zielsetzung}

Hier formulieren Sie die konkreten Ziele Ihrer Arbeit. Sie können auch auf andere Kapitel verweisen, \zB auf \autoref{chap:grundlagen} für die theoretischen Grundlagen oder auf \autoref{chap:schluss} für das Fazit.

\subsection{Hauptziel}
Das Hauptziel dieser Arbeit ist \dots

\subsection{Teilziele}
Die Arbeit verfolgt folgende Teilziele:
\begin{enumerate}
    \item Analyse von \dots
    \item Entwicklung von \dots
    \item Evaluation von \dots
\end{enumerate}

\section{Forschungsmethodik}
\label{sec:methodik}

Beschreiben Sie hier den methodischen Ansatz Ihrer Arbeit. Literaturverweise werden mit \texttt{\textbackslash cite} eingefügt, etwa so: Die verwendete Methodik basiert auf etablierten Ansätzen aus der Literatur.

Sie können auch Glossarbegriffe verwenden: \gls{egov} ist ein wichtiges Konzept in der digitalen Verwaltung. Beim ersten Verwenden eines Akronyms wird es automatisch ausgeschrieben: \gls{desi}.

\section{Aufbau der Arbeit}
\label{sec:aufbau}

Die Arbeit ist wie folgt strukturiert:

\begin{description}
    \item[\autoref{chap:grundlagen}] behandelt die theoretischen Grundlagen und definiert wichtige Begriffe.
    \item[\autoref{chap:hauptteil}] beschreibt die Hauptuntersuchung bzw. Implementierung.
    \item[\autoref{chap:schluss}] fasst die Ergebnisse zusammen und gibt einen Ausblick auf zukünftige Arbeiten.
\end{description}

% Beispiel für Abbildungsreferenz (anpassen falls Sie Abbildungen haben)
% In \autoref{fig:beispiel} ist die Systemarchitektur dargestellt.

% Beispiel für Tabellenreferenz
% Die Anforderungen sind in \autoref{tab:anforderungen} zusammengefasst.

\section{Abgrenzung}
\label{sec:abgrenzung}

Grenzen Sie hier ab, was \textbf{nicht} Teil Ihrer Arbeit ist:
\begin{itemize}
    \item Aspekt A wird nicht behandelt, da \dots
    \item Aspekt B liegt außerhalb des Projekt-Scopes
    \item Aspekt C wurde aus zeitlichen Gründen nicht berücksichtigt
\end{itemize}

% ===================================================================
% HINWEIS: Passen Sie alle Platzhalter und Beispieltexte an!
% ===================================================================